% DRAFT (Claude, 2026-11-10, Anders): 5.3 written with the literature-review skill in the compact
% form agreed for 5.2: a short opener, one or two paragraphs per sub-question, a close.
% Every citation's passage in CHANGES-5-3.md ("note" where only the paper note was checked).
% No numbers: the comparability caveat comes in 5.4. Second version: sentences split, rhythm as in the theses.

% Models (precedent rule, SKILL.md hard rule 7): section layout = Bore 2026 §4.2 (opener, four numbered sub-sections, synthesis).
%   5.3 opener        <- bore-2026 P04 (relevance with a pointer to the background, then what the subsection covers)
%   5.3.1 paragraph 1 <- bore-2026 P05 + bore-brensholm-2025 P04 (domain fact, then studies as evidence, dimensionality limit as in bore-2026 P07)
%   5.3.1 paragraph 2 <- bore-2026 P12 (correlation: two studies contrasted, consequence at the end)
%   5.3.2             <- bore-2026 P14 (granularity: claim, two or three studies contrasted, "should be chosen in line with the purpose")
%   5.3.3             <- bore-2026 P13 (scenario-tree steps shared by the studies, data length and granularity, consequence)
%   5.3.4             <- bore-2026 P06 (domain fact with sources) + norland-2025 P04 (group of studies, "yet none assess …")
%   close             <- bore-2026 P10 (synthesis: the key variables, the number kept small, what is left)

\subsection{Uncertainty in Pumped-Storage Scheduling}
\label{sec:lit_uncertainty}

Uncertainty about future prices and inflow, introduced in
Section~
ef{sec:background-uncertainty}, is central to hydropower
scheduling, and more so for pumped storage, where prices set both the revenue
from generating and the cost of pumping.
% TODO: label of the background section on uncertainty
The reviewed studies differ in which inputs they treat as uncertain, when the
information is assumed known, how the uncertainty is represented over time,
and how the behaviour of prices is described.

\subsubsection{Sources of Uncertainty and Their Dependence}
\label{sec:lit_unc_sources}

Stochastic modelling can give a better representation of reality, since
inflow, prices and weather are not known when the water is allocated.
\citet{belsnes2016applying} find that a stochastic strategy earns more than a
deterministic one in every case they test. The main reason is that it sees a
risk of spill that the deterministic plan does not.
% src: belsnes2016applying pp. 176-178, Tables 3-5; p. 177
Similarly, \citet{lohndorf2013} find for a pumped-storage plant that the gain
over a rolling deterministic plan grows with the power capacity relative to
the reservoir. A plant that can move its reservoir content quickly benefits
more from anticipating prices.
% src: lohndorf2013 p. 28
Inflow and energy prices are the inputs most studies treat as stochastic
\citep{lohndorf2013,helseth2016,kleiven2023,alic2023optimal,bore2026assessing}.
Some also treat reserve-capacity prices as stochastic
\citep{helseth2016,helseth2017iet}.
% src: notes of each; helseth2016 p. 935; helseth2017iet p. 130 ("treating both prices as stochastic")
Ideally every price the plant trades on would be stochastic. However, each
added variable multiplies the number of states and the data needed to
estimate the transitions between them. \citet{helseth2016}, for example,
have to adjust their transition probabilities because few scenarios are
available.
% src: helseth2016 pp. 938-939 (note: 3x3 nodes from 80 scenarios, probabilities adjusted)

When several inputs are uncertain, their dependence is a further choice.
\citet{kleiven2023} start from the negative correlation between Norwegian
prices and inflow. Ignoring it changes water values and reservoir operation
in their case, although the loss in expected revenue was small.
% src: kleiven2023 p. 303 (premise), p. 524 (0.06 %), takeaway (0.17-0.30 %)
In contrast, \citet{helseth2016} assume prices independent of inflow because
their decomposition requires it.
% src: helseth2016 p. 937
Both concern conventional plants. How dependence matters when prices also
drive pumping is therefore an open question, and an important one for the
modelling assumptions in this thesis.

\subsubsection{Timing and Resolution of Information}
\label{sec:lit_unc_timing}

Most medium-term studies assume that the current week's prices and inflow are
known once the week starts, and keep the uncertainty between weeks
\citep{helseth2013,bore2026assessing}.
% src: helseth2013 p. 1432 ("values of the stochastic variables are known for a whole week in advance"); bore2026assessing p. 3 ("inflow for the period are assumed to be known")
\citet{helseth2013} note that this is too coarse to capture the need for
balancing wind power.
% src: helseth2013 p. 1432 ("a week-long stochastic time resolution is not fine enough to capture the actual need for balancing wind power")
\citet{abgottspon2014} test the assumption for a Swiss pumped-storage plant.
Treating the week as known overvalues the water, whereas a stochastic
within-week model gives a higher expected profit in their case with weekly
reserves.
% src: abgottspon2014 p. 6 ("it overvalues the water value"), Table II (note: about 10 % with reserves, not without)
Resolution is a separate choice from information. \citet{ak2019quantifying}
find that operating their cascade as pumped storage barely changes the
monthly reservoir decisions, but changes the hourly allocation of pumping and
generation considerably.
% src: ak2019quantifying pp. 746-747 (note)
Consequently, a long horizon alone does not show that a model captures the
flexibility a pump adds.

\subsubsection{Representing Uncertainty over Time}
\label{sec:lit_unc_repr}

The reviewed studies typically represent uncertainty by a few states per
week. The states are built by clustering historical or simulated series, and
the transitions between them are counted
\citep{helseth2016,alic2023optimal,bore2026assessing}.
% src: helseth2016 pp. 938-939; alic2023optimal p. 25 (five states); bore2026assessing pp. 3-4 (three states)
More states describe the variation better but rest on fewer observations
each. \citet{mo2001} show how the transition probabilities can be corrected
so that the conditional expected prices are preserved.
% src: mo2001 (note)
The studies also differ in what a state contains.
\citet{alic2023optimal} give each state a weekly average price with a fixed
daily profile, and \citet{ak2019quantifying} use representative daily curves
per month. \citet{bore2026assessing}, in contrast, keep the hourly series of
each cluster.
% src: alic2023optimal text l. 196, 602 (weekly average energy price), note (fixed profile); ak2019quantifying p. 741 (note); bore2026assessing p. 4 (note: median series)
For a pump that earns from differences within the week, averaging the prices
inside a state may understate its value, although none of the studies tests
this.

\subsubsection{Price Volatility and Persistence}
\label{sec:lit_unc_price}

Nordic day-ahead prices show long memory, strongest in the Norwegian areas,
which \citet{haldrup2006regime} relate to the slow variation of reservoirs.
For a storage owner this behaviour over time matters beyond which inputs are
uncertain.
% src: haldrup2006regime abstract (long memory), note (Norwegian areas, 2000-2003 data)
A weekly Markov chain captures such persistence only through its transition
probabilities. Moreover, it forgets its initial state at an exponential rate
\citep{mo2001}.
% src: mo2001 text l. 295-296 ("a Markov model will tend to forget its initial state at an exponential rate")
Price behaviour also changes with the system. \citet{henden2016economic}
find that more transmission to the continent raises the variability of
southern Norwegian prices, whereas more pumped storage lowers it.
% src: henden2016economic p. 120, Fig. 1a
Studies that relate price behaviour to storage value use known prices.
\citet{lohndorf2023} find that coordinating day-ahead and intraday trading
benefits a pumped-storage plant far more than a reservoir plant, and
\citet{ak2019quantifying} earn most from pumping in the year with the largest
daily price variation.
% src: lohndorf2023 text l. 1330-1345 (pumped hydro +10.0-29.1 % over sequential; reservoir gains marginal); ak2019quantifying p. 747 (note)
None of them, however, varies volatility and persistence separately or relates
them to the value of pump, turbine and reservoir capacity, which is the
question this study asks.

Overall, the reviewed studies treat prices and inflow as the uncertain
inputs, keep their number small, assume the current week known, and build
weekly states from data. What they leave open for pumped storage is how the
price pattern inside a week and the persistence of price conditions between
weeks affect the value of each capacity.
